\section{Related Work}

\noindent\textbf{Learned and heuristic placement for log-structured storage.}
\sepbit, \midas, and \dogi improve log-structured placement by predicting or adapting to invalidation behavior~\cite{sepbit,midas,dogi}.  \sepbit separates blocks by inferred invalidation time, \midas adapts group sizes and counts to workload dynamics, and \dogi uses oracle-guided analysis to motivate hot filtering, learned lifetime prediction, GC relocation policy, and dynamic group organization.  These systems are the right comparison class for \system because they represent the strongest storage-history view of placement.

\system is not an argument that learned placement is generally unnecessary.  For ordinary payload and workloads whose lifetime is not visible to the application, history is the right signal and \dogi-style placement can outperform semantic grouping.  \system targets a narrower regime where lifetime is not hidden from the software stack: a KEM artifact, ephemeral secret, or KMS envelope dies because a protocol epoch closes.  In that regime, prediction spends work approximating a label that the application already has, and a misplacement can leave expired secret bytes in a long-lived zone.

\noindent\textbf{Zoned storage and host-guided placement.}
\zns and related host-managed SSD interfaces expose placement and reclaim decisions to the host~\cite{zns,zenfs,xnvme}.  Systems such as ZenFS show how a host-side stack can map application writes onto zone append and reset operations.  \system builds on this host-managed direction but changes the placement key: instead of grouping only by hotness, stream, or predicted invalidation time, it groups secret-bearing writes by protocol death cohort.

Flexible Data Placement provides a different host-guided interface without forcing a full \zns software stack~\cite{fdp}.  \system can target native \zns by choosing zone families directly, or FDP-like devices by mapping intent and epoch to placement handles.  This makes \system a semantic placement policy rather than a device-specific API proposal.

\noindent\textbf{Benchmarks and trace methodology.}
FAST storage papers typically evaluate both synthetic controls and macro workloads.  DOGI uses FIO Zipf, YCSB-A/F, Varmail, Alibaba, and Exchange-like traces; our workload plan preserves these axes and adds PQC lifecycle overlays.  FIO, YCSB, and Sysbench remain important because they provide recognizable hot/cold, update-heavy, and DB-style pressure points~\cite{fio,ycsb,sysbench}.  The paper therefore treats pure PQC epoch streams only as sanity checks and uses DOGI/FAST-shaped pressure rows for headline claims.

\noindent\textbf{Post-quantum systems.}
Open Quantum Safe, liboqs, and oqs-provider provide software paths for deploying post-quantum algorithms in applications and TLS stacks~\cite{liboqs,oqsprovider}.  NIST's ML-KEM and ML-DSA standards make these deployments concrete targets rather than hypothetical future algorithms~\cite{nist-fips203,nist-fips204}.  Most systems work around PQC focuses on cryptographic correctness, interoperability, CPU overhead, and network handshake cost.  \system instead studies a storage side effect: lifecycle-heavy cryptographic metadata creates a placement signal that the block layer and device do not naturally observe.

\noindent\textbf{Secure erase and sanitization.}
Storage sanitization and crypto-erase mechanisms provide device-specific ways to destroy data or keys~\cite{nvme-sanitize}.  \system does not replace these mechanisms and does not claim that zone reset alone physically erases NAND.  It solves a placement precondition for secure disposal: expired secrets are isolated into zone families that can be reset or sanitized at the protocol boundary.  When hardware provides crypto-erase or sanitize semantics, \system gives the storage stack a deterministic point at which to trigger them.
